\documentclass[11pt]{article}
\usepackage[margin=1in]{geometry}
\usepackage{amsmath,amssymb,amsthm,mathtools}
\usepackage{microtype}
\usepackage[hidelinks]{hyperref}

\newtheorem{theorem}{Theorem}
\newtheorem{lemma}{Lemma}
\newcommand{\F}{\mathcal F}
\newcommand{\E}{\mathbb E}
\newcommand{\Prb}{\mathbb P}

\title{An unconditional $0.3827090879$ lower bound for the union-closed sets conjecture}
\author{Unsolved Labs Research Release R010}
\date{14 August 2026}

\begin{document}
\maketitle

\begin{abstract}
We give a computer-assisted proof that every finite union-closed family
$\F\neq\{\varnothing\}$ contains an element belonging to at least
\[
  \frac{3827090879}{10^{10}}|\F|=0.3827090879|\F|
\]
members of the family. The argument uses the conditionally-i.i.d. coupling
introduced by Liu, but removes the numerical structural hypotheses used for
Liu's approximately $0.38271$ evaluation. The proof reduces the union-closed
statement to two pointwise entropy inequalities. One is proved analytically,
with one polynomial positivity step checked exactly over $\mathbb Q$; the
other is proved by analytic boundary estimates and a rigorous two-dimensional
interval certificate. The numerical component is replayed independently with
directed-rounding Boost.Interval arithmetic, a separately written 256-bit MPFR
implementation with an independently generated subdivision, and a non-adaptive
MPFR replay of a frozen complete prefix-coded partition. This research release
was generated with frontier AI; independent specialist review remains pending.
\end{abstract}

\section{Introduction}

The union-closed sets conjecture, usually attributed to Frankl, asks whether
every nontrivial finite union-closed family contains an element present in at
least half of its member sets. Gilmer obtained the first absolute constant
lower bound by an entropy argument. Sawin, Yu, and Cambie subsequently
developed dependent-coupling refinements reaching an unconditional frontier
near $0.3823455$. Liu introduced conditionally-i.i.d. couplings and proved an
analytic strict improvement. A nine-dimensional numerical optimization in
Liu's paper suggested a value near $0.38271$ under explicitly numerically
verified structural hypotheses.

R010 replaces those hypotheses, for the exact rational constant below, by a
direct two-variable inequality with a rigorous interval certificate.

\begin{theorem}[R010]\label{thm:main}
Let $\F\neq\{\varnothing\}$ be a finite union-closed family. Then some
element belongs to at least
\[
c|\F|,\qquad c=\frac{3827090879}{10000000000}=0.3827090879
\]
members of $\F$.
\end{theorem}

The number is deliberately slightly below Liu's numerically optimized value.
That slack makes the decisive compact inequality strictly positive and permits
a modest exact interval certificate. The theorem does not prove the $1/2$
conjecture, and no optimality is claimed for $c$ or for this proof architecture.

\section{Notation and exact constants}

Write
\[
h(t)=-t\log t-(1-t)\log(1-t),\qquad 0\le t\le1,
\]
with natural logarithms, and define
\[
J(x,y)=xy\bigl(1+(1-x)(1-y)\bigr),\qquad
j(x)=J(x,x)=x^2\bigl(1+(1-x)^2\bigr).
\]
Use the exact constants
\[
\alpha=\frac{11249343}{12500000},\quad
\beta=\frac{1250657}{12500000},\quad
c=\frac{3827090879}{10000000000},
\]
\[
m=1-c=\frac{6172909121}{10000000000},\qquad
\kappa=\frac1{2m}=\frac{5000000000}{6172909121},
\]
so $\alpha+\beta=1$.

\section{Two pointwise inequalities}

The proof is driven by two statements.

\begin{lemma}[Kernel inequality]\label{lem:kernel}
For every $x,y\in[0,1]$,
\[
h(J(x,y))\ge\sqrt{h(j(x))h(j(y))}.
\]
\end{lemma}

\begin{lemma}[Main entropy inequality]\label{lem:mainineq}
For every $x,y\in[0,1]$,
\[
\alpha h(xy)+\beta\sqrt{h(j(x))h(j(y))}
\ge\kappa\bigl(xh(y)+yh(x)\bigr).
\]
\end{lemma}

The kernel inequality is analytic apart from an exact rational polynomial
positivity check. The main inequality is analytic on two boundary strips and
rigorously interval-certified on the remaining compact rectangle.

\section{Deduction of the union-closed bound}\label{sec:deduction}

Let $Z$ be uniformly distributed on the indicator vectors of $\F$. The case
$|\F|=1$ is immediate, so assume $|\F|>1$. For a prefix $z^{i-1}$ write
\[
p_i(z^{i-1})=\Prb[Z_i=1\mid Z^{i-1}=z^{i-1}].
\]
Use two sequential couplings of copies of $Z$. Under $\Pi^{(0)}$, the two
coordinate bits are independent Bernoulli variables with the prescribed
conditional means. Under $\Pi^{(1)}$, use Liu's conditionally-i.i.d. protocol
(Example 5 in \cite{Liu2023}) with $f(r)=r(1-r)$. Let
$(A^{(k)},B^{(k)})$ be the pair generated by protocol $k$ and set
$W^{(k)}=A^{(k)}\vee B^{(k)}$. Each marginal has the law of $Z$.

Fix coordinate $i$. If
\[
S=p_i(A^{(k),i-1}),\qquad T=p_i(B^{(k),i-1}),
\]
then $S,T$ have a protocol-independent common marginal law. Let $P_i$ be the
law of $X=1-S$, and put
\[
M_i=\E X=1-\Prb[Z_i=1],\qquad
H_i=\E h(X)=H(Z_i\mid Z^{i-1}).
\]
For the independent protocol, $X,Y$ are independent with law $P_i$, and the
conditional probability that the OR bit is zero is $XY$. Integrating
Lemma~\ref{lem:mainineq} gives
\[
\alpha\E h(XY)+\beta\bigl(\E\sqrt{h(j(X))}\bigr)^2
\ge\frac{M_i}{m}H_i. \tag{4.1}
\]
For the conditionally-i.i.d. protocol, Liu's construction gives OR-zero
probability $J(X,Y)$. Conditionally on its auxiliary variable $U$, $X,Y$ are
i.i.d. Writing $s(x)=\sqrt{h(j(x))}$, Lemma~\ref{lem:kernel}, conditional
independence, and Jensen give
\[
\E h(J(X,Y))\ge\E[s(X)s(Y)]
=\E[\E(s(X)\mid U)^2]\ge(\E s(X))^2. \tag{4.2}
\]
Thus
\[
\alpha\E h(XY)+\beta\E h(J(X,Y))\ge\frac{M_i}{m}H_i. \tag{4.3}
\]
The chain rule and the fact that $W^{(k),i-1}$ is a function of the paired
histories imply
\[
H(W^{(k)})\ge\sum_i H(W_i^{(k)}\mid A^{(k),i-1},B^{(k),i-1}). \tag{4.4}
\]
If every element occurred in strictly fewer than $c|\F|$ members, then
$M_i>m$ for every coordinate. Summing (4.3) and using (4.4) yields
\[
\alpha H(W^{(0)})+\beta H(W^{(1)})
>\sum_iH_i=H(Z)=\log|\F|.
\]
The inequality is strict because $H(Z)>0$, hence some $H_i>0$. Union closure
implies each $W^{(k)}$ is supported on $\F$, so each entropy is at most
$\log|\F|$, a contradiction.

\section{Proof of the kernel inequality}\label{sec:kernel}

Boundary cases follow by continuity. Assume $0<x\le y<1$ and define
\[
\ell(t)=\log h(t),\qquad
R(x,y)=2\ell(J(x,y))-\ell(j(x))-\ell(j(y)).
\]
Since $h''<0$ and $h>0$, $\ell$ is strictly concave, while $R(x,x)=0$.
It suffices to prove $R_y(x,y)\ge0$. Direct expansion gives
\[
J(x,y)\le j(y),\qquad 2J_y(x,y)\le j'(y),\qquad x\le y. \tag{5.1}
\]
When $j(y)\ge1/2$, the sign and monotonicity of $\ell'$ together with (5.1)
imply $R_y\ge0$.

Assume now $j(y)<1/2$. Put $G(x,y)=\ell(J(x,y))$. Since
$R_y=2(G_y(x,y)-G_y(y,y))$, it is enough to show $G_{xy}\le0$. With
$z=J(x,y)$,
\[
G_{xy}=\ell''(z)J_xJ_y+\ell'(z)J_{xy},
\]
and it suffices to establish
\[
-\frac{z\ell''(z)}{\ell'(z)}\ge\frac1{1-z}
\ge\frac{zJ_{xy}}{J_xJ_y}. \tag{5.2}
\]
For the entropy side let $a=\log((1-z)/z)>0$. Direct calculation gives
\[
-\frac{z\ell''(z)}{\ell'(z)}
=\frac1{(1-z)a}+\frac{za}{h(z)}.
\]
After multiplying by $h(z)(1-z)a$, the required difference has numerator
$N=h(z)(1-a)+za^2(1-z)$. The case $0<a\le1$ is immediate. For $a\ge1$ use
\[
h(z)=za-\log(1-z)\le z\left(a+\frac1{1-z}\right).
\]
Because $1-a\le0$, it remains to use $z<(e^a)^{-1}$ and
$e^a\ge a+a^2$, the latter following from the Taylor expansion of $e^a$.

For the geometric side exact algebra gives
\[
J_xJ_y-(1-J)JJ_{xy}=x^2y^2P(x,y),
\]
where
\begin{align*}
P(x,y)={}&4x^3y^3-10x^3y^2+8x^3y-2x^3
-10x^2y^3+34x^2y^2-34x^2y+10x^2\\
&+8xy^3-34xy^2+44xy-16x-2y^3+10y^2-16y+7.
\end{align*}
Moreover
\[
j\left(\frac{84}{125}\right)-\frac12
=\frac{81647}{488281250}>0.
\]
Since $j$ is increasing, $j(y)\le1/2$ implies $x,y\le84/125$. After
rescaling to the unit square, all sixteen tensor-product Bernstein
coefficients of $P$ are positive. The exact checker
\texttt{verify\_auxiliary\_inequality.py} recomputes the polynomial identity,
the endpoint comparison, and all sixteen rational Bernstein coefficients.
Hence $P>0$, proving (5.2) and Lemma~\ref{lem:kernel}.

\section{Proof of the main entropy inequality}\label{sec:mainineq}

For $t\in(0,1)$ define
\[
u(t)=\frac{h(t)}t,\qquad
v(t)=\frac{\sqrt{h(j(t))}}t.
\]
After division by $xy$, Lemma~\ref{lem:mainineq} is equivalent to
\[
\Delta(x,y)=\alpha u(xy)+\beta v(x)v(y)
-\kappa\bigl(u(x)+u(y)\bigr)\ge0. \tag{6.1}
\]
The square is split into a small-variable strip $\min(x,y)\le0.01$, the
compact core $[0.01,0.999999]^2$, and a near-one strip
$\max(x,y)\ge0.999999$.

\subsection{Small-variable strip}
For $0<x,y\le1$,
\[
u(xy)\ge u(x)-\log y, \tag{6.2}
\]
obtained by integrating $g'(t)=-\log(1-e^{-t})/e^{-t}\ge1$ for
$g(t)=u(e^{-t})$. For $x\le0.01$,
\[
j(x)\ge\frac{19801}{10000}x^2,\qquad
j(x)\le2x^2\le\frac1{5000},
\]
so $h(z)\ge z\log(1/z)$ and $\log5000>8$ imply $v(x)>7/2$. Since
$\beta>1/10$, it remains to prove positivity of
\[
R(y)=-\alpha\log y-\kappa u(y)+\frac7{20}v(y).
\]
Elementary entropy estimates handle $0<y\le10^{-4}$ and
$0.9999\le y<1$. On $[10^{-4},0.9999]$, the rigorous interval routine
\texttt{verify\_R()} uses 39 boxes and has reference least lower endpoint
$0.0033854057973891488$.

\subsection{Near-one strip}
By symmetry assume $x\le y$ and $y\ge0.999999$, with $x>0.01$. If
$x\le0.998$, monotonicity of $u$ gives a directed-rounding lower bound
$0.0012884232397950039$. For $0.998\le x\le y<1$, concavity of
$g(t)=u(e^{-t})$ gives
\[
u(xy)\ge\frac{u(x^2)+u(y^2)}2.
\]
It is therefore enough to show $\alpha u(z^2)>2\kappa u(z)$ for
$0.998\le z<1$. Writing $z=1-\varepsilon$ and using
\[
h(t)=t(1-\log t)-E(t),\qquad
0\le E(t)\le\frac{t^2}{2(1-t)},
\]
reduces the claim to explicit rational/logarithmic estimates; the resulting
published lower bound is $391/12500>0$.

\subsection{Compact core}
On $[0.01,0.999999]^2$, the primary checker uses Boost.Interval with IEEE-754
directed rounding. It does not delegate logarithms to system libm. For
$1\le q\le2$, with $w=(q-1)/(q+1)$, it encloses
\[
\log q=2\sum_{n=0}^{24}\frac{w^{2n+1}}{2n+1}+R_{24}
\]
using the explicit positive remainder bound
\[
0\le R_N\le\frac{2w^{2N+3}}{(2N+3)(1-w^2)}.
\]
The checker combines natural interval evaluation with a second-order Taylor
enclosure on each rectangle and bisects only when no current enclosure proves
a nonnegative lower endpoint. The reference Boost run reports 3273 accepted
core boxes, maximum depth 39, and least accepted lower endpoint
$4.1293562138181227\times10^{-12}$.

\section{Independent numerical verification}

The numerical obligation is deliberately replayed through distinct trust
paths. The Boost path uses directed binary64 arithmetic and an internally
bounded logarithm series. A separately written MPFR checker uses 256-bit
outward rounding, MPFR's own logarithm and square-root functions, exact
rational domain endpoints, and an independently generated subdivision; its
reference core contains 3210 boxes and has least lower endpoint
$4.1428237777820989\times10^{-12}$. The same MPFR subdivision is frozen as
prefix-coded leaf paths. A non-adaptive static checker validates the trie as a
complete partition before checking the 3210 leaves. Finally, the Bernstein
step in Section~\ref{sec:kernel} is recomputed with exact
\texttt{fractions.Fraction} arithmetic.

Search and heuristic optimization are therefore outside the final correctness
oracle. The theorem relies on the analytic reductions above and on the exact
or directed-rounding certificate checks just described.

\section{Reproducibility and trust boundary}

From a clean Ubuntu/Debian checkout install \texttt{g++}, \texttt{python3},
Boost, GMP, and MPFR and run
\begin{verbatim}
./verify_release.sh
\end{verbatim}
The release materializer validates the SHA-256 of the compressed verifier and
certificate payload and rejects unsafe archive paths and links. CI replays the
exact auxiliary check, Boost certificate, independent MPFR certificate,
static MPFR certificate, and public release-integrity checks.

No proof-assistant formalization is currently part of the R010 trust boundary.
A future Lean formalization would be most valuable for the protocol-to-entropy
deduction and the analytic reduction from the two pointwise inequalities to
Theorem~\ref{thm:main}. The large interval partition is already represented as
a frozen deterministic certificate checked by independent rigorous programs.

\section{Relation to prior work and novelty boundary}

Gilmer \cite{Gilmer2022} initiated the constant entropy method. Sawin
\cite{Sawin2022}, Yu \cite{Yu2022}, and Cambie \cite{Cambie2025} developed
dependent-coupling refinements, with an unconditional frontier around
$0.3823455$. Liu \cite{Liu2023} introduced the conditionally-i.i.d. protocol
used here and proved an analytic strict improvement. Liu's numerical section
reported an approximately $0.38271$ value under numerically verified
hypotheses arising from a nine-dimensional optimization.

R010 does not claim invention of the protocol or of the numerically suggested
target. Its contribution is the unconditional exact rational constant
$0.3827090879$, obtained by replacing those structural/global-optimization
hypotheses with the two pointwise inequalities proved and rigorously certified
here.

\section{Limitations}

The full Frankl $1/2$ conjecture remains open. This work does not establish the
best constant obtainable by entropy methods, and it does not prove Liu's full
numerically suggested decimal $0.382709087918741\ldots$. Independent
specialist review remains pending.

\begin{thebibliography}{9}
\bibitem{Gilmer2022} J. Gilmer, \emph{A constant lower bound for the union-closed sets conjecture}, arXiv:2211.09055, 2022.
\bibitem{Sawin2022} W. Sawin, \emph{An improved lower bound for the union-closed set conjecture}, arXiv:2211.11504, 2022.
\bibitem{Yu2022} L. Yu, \emph{Dimension-Free Bounds for the Union-Closed Sets Conjecture}, arXiv:2212.00658, 2022.
\bibitem{Cambie2025} S. Cambie, \emph{Better bounds for the union-closed sets conjecture using the entropy approach}, arXiv:2212.12500v2, 2025 revision.
\bibitem{Liu2023} J. Liu, \emph{Improving the Lower Bound for the Union-closed Sets Conjecture via Conditionally IID Coupling}, arXiv:2306.08824v1, 2023.
\end{thebibliography}

\end{document}
